\begin{table}[htbp]
\centering
\caption{Robustness: Placebo Industry and Leave-One-Out Tests}
\label{tab:robustness}
\begin{tabular}{lccc}
\hline\hline
 & DDD Coefficient & SE & Observations \\
\hline
\multicolumn{4}{l}{\textit{Panel A: Placebo industry (NAICS 23 --- Construction)}} \\[3pt]
ln(Employment) & 0.1320** & (0.0548) & 147,898 \\
ln(Earnings) & 0.0588** & (0.0225) & 147,551 \\
\hline
\multicolumn{4}{l}{\textit{Panel B: Leave-one-out (ln Employment, NAICS 72)}} \\[3pt]
Drop Oregon & -0.0777** & (0.0344) & 161,180 \\
Drop San Francisco & 0.0876* & (0.0481) & 163,043 \\
Drop New York City & 0.1202*** & (0.0164) & 162,811 \\
\hline
\multicolumn{4}{l}{\textit{Panel C: Full sample including COVID-era adopters (2013Q1--2022Q4)}} \\[3pt]
ln(Employment) & 0.1329* & (0.0670) & 224,804 \\
ln(Earnings) & 0.0113** & (0.0045) & 224,684 \\
\hline\hline
\end{tabular}
\begin{tablenotes}
\small
\item \textit{Notes:} Panel A shows triple-difference estimates for NAICS 23 (Construction), an industry 
not covered by Fair Workweek laws but with substantial Black employment. Panel B drops each treatment 
jurisdiction in turn from the pre-COVID NAICS 72 sample. Panel C extends the sample through 2022Q4, 
adding Philadelphia (2020) and Chicago (2020) as treated jurisdictions. All specifications include 
county$\times$race, quarter$\times$race, and state$\times$quarter FE. SEs clustered at state level. 
Randomization inference p-value (500 permutations): 0.606.
 $^{***}p<0.01$, $^{**}p<0.05$, $^{*}p<0.1$.
\end{tablenotes}
\end{table}
